\Needspace{17\baselineskip}
\SourceLeaf{134}
\noindent We further note that if the numbers in the first column are regarded as twice the logarithms of tangents, the numbers in the second and third columns will be twice the complements of the logarithms of the corresponding sines and cosines. Ordinary trigonometric tables can therefore serve the intended purpose in exactly the same way as those proposed by \textsc{Leonelli}. The latter do save division and multiplication by 2, which anyone with moderate practice in calculation can easily do mentally, and yield twice the accuracy when calculated to the same number of decimal places. It is curious that \textsc{Delambre} does not mention the method of using trigonometric tables for the same purpose based on this observation, since he describes another, far less convenient one. \textsc{Leonelli} is quite right to complain that anyone should propose putting such a procedure on a par with his own.\par\medskip\begin{center}\rule{.24\linewidth}{.3pt}\end{center}

